% we use biblatex, which is a more modern alternative to the traditional bibtex
% comment this to use bibtex instead
\usepackage[
    backend=biber,
    style=alphabetic,
    sorting=nyt,
    maxbibnames=99
]{biblatex}

% graphicx is needed to include images
\usepackage{graphicx}
\usepackage{adjustbox}
\usepackage{float}
\usepackage{flafter} % a figure never appears before the text that introduces it
% let large figures share a page with text and with each other (less empty space)
\renewcommand{\topfraction}{0.85}
\renewcommand{\bottomfraction}{0.7}
\renewcommand{\textfraction}{0.1}
\renewcommand{\floatpagefraction}{0.75}

% xspace is used to add a space after a macro, if needed
\usepackage{xspace}

% acronym (https://ctan.org/pkg/acronym) is a package to expand acronyms after their first usage
\usepackage{acronym}

% the hyperref package provides links; with this configuration they are colored
% according to the scheme described in https://tex.stackexchange.com/a/525297
\usepackage{xcolor}
\usepackage[colorlinks]{hyperref}
\def\tmp#1#2#3{%
  \definecolor{Hy#1color}{#2}{#3}%
  \hypersetup{#1color=Hy#1color}}
\tmp{link}{HTML}{800006}
\tmp{cite}{HTML}{2E7E2A}
\tmp{file}{HTML}{131877}
\tmp{url} {HTML}{8A0087}
\tmp{menu}{HTML}{727500}
\tmp{run} {HTML}{137776}
\def\tmp#1#2{%
  \colorlet{Hy#1bordercolor}{Hy#1color#2}%
  \hypersetup{#1bordercolor=Hy#1bordercolor}}
\tmp{link}{!60!white}
\tmp{cite}{!60!white}
\tmp{file}{!60!white}
\tmp{url} {!60!white}
\tmp{menu}{!60!white}
\tmp{run} {!60!white}

% the cleveref package is used to automatically add the type of reference
% (e.g., "section", "figure", "table") to the reference using the \cref command
% The "capitalize" option makes both \cref and \Cref uppercase; if we remove it
% only \Cref will be uppercase.
\usepackage[capitalize]{cleveref}
% cleveref loads before listings below, so it does not know what to call an
% \lstlisting by default; name it explicitly or \Cref{lst:...} prints "??"
\crefname{lstlisting}{Listing}{Listings}
\Crefname{lstlisting}{Listing}{Listings}

% the booktabs package is used to create better tables; check an example of usage in results.tex
\usepackage{booktabs}
\usepackage{tabularx}

% breaks long URLs better
\usepackage{xurl}

% macros
\newcommand*{\latex}{\LaTeX \xspace}
\newcommand{\myhref}[2]{#2\footnote{\url{#1}}}

% tracking changes using the todonotes package
\setlength{\marginparwidth}{2cm}
\usepackage{todonotes}
\newcommand{\note}[3]{\todo[inline,color=#2]{#1: #3}}
\newcommand{\md}[1]{\note{MD}{yellow}{#1}}

% result is something like "Related Work (Chapter 2)"
\newcommand{\fullref}[1]{\nameref{#1} (\Cref{#1})}

% used for the inparaenum environment
\usepackage{paralist}

% code listings for chaincode and API examples
\usepackage{listings}
\usepackage{lstautogobble}
\lstset{
    basicstyle=\ttfamily\small,
    breaklines=true,
    frame=single,
    numbers=left
}

% for nicer diagrams and architecture figures
\usepackage{tikz}
\usetikzlibrary{arrows.meta,positioning,fit,calc}

% styles shared by the TikZ figures
\tikzset{
    box/.style={draw, rounded corners=2pt, align=center, minimum height=9mm, inner sep=4pt, fill=gray!8},
    chain/.style={draw, rounded corners=2pt, align=center, minimum height=9mm, minimum width=22mm, inner sep=3pt, fill=blue!6},
    arr/.style={-{Latex[length=2mm]}, thick},
    darr/.style={-{Latex[length=2mm]}, thick, dashed},
    lbl/.style={font=\scriptsize, midway, fill=white, inner sep=1pt}
}

% \screenshot{file}{caption}{label}: includes images/<file> when it exists,
% otherwise draws a labelled placeholder so the document still compiles.
% Use file names without underscores.
% Screenshots never appear before the text that introduces them (flafter). \screenshot uses the full text width; \modalshot crops
% the blurred background around a window (about 23.5% left/right, 4% top/bottom)
% so that the window itself fills the width.
\newcommand{\modalshot}[3]{%
    \begin{figure}[!htbp]
        \centering
        \IfFileExists{#1}%
            {\adjustbox{trim={.235\width} {.04\height} {.235\width} {.04\height},clip,width=0.66\textwidth}{\includegraphics{#1}}}%
            {\fbox{\parbox{0.8\textwidth}{\centering\vspace{2.2cm}\texttt{#1}\\[2pt]\small screenshot to be added\vspace{2.2cm}}}}
        \caption{#2}
        \label{#3}
    \end{figure}}
\newcommand{\screenshot}[3]{%
    \begin{figure}[!htbp]
        \centering
        \IfFileExists{#1}%
            {\includegraphics[width=\textwidth]{#1}}%
            {\fbox{\parbox{0.8\textwidth}{\centering\vspace{2.2cm}\texttt{#1}\\[2pt]\small screenshot to be added\vspace{2.2cm}}}}
        \caption{#2}
        \label{#3}
    \end{figure}}

% Anita's own notes
\newcommand{\ag}[1]{\note{AG}{cyan}{#1}}